\section{Conclusion}

We asked whether data attribution method performance depends on model architecture, and found that it does---but the patterns differ from theoretical predictions.

\textbf{TRAK exhibits remarkable architecture-invariance} ($<$1\% AUC difference). \textbf{TracIn shows a directional BERT advantage} ($\sim$3\%). \textbf{EK-FAC does not show the predicted GPT-2 advantage} (0.27\%).

We validated the causal mechanism: 98.8\% attention sparsity difference creates 11$\times$ Hessian eigenvalue difference, propagating through each method's approximation.

Returning to our opening question: when choosing attribution methods across architectures, practitioners can now rely on TRAK for consistent performance. Architecture-agnostic attribution is possible---random projection provides the key.

\subsection{Future Work}

\begin{itemize}
    \item Scale testing (1B+ parameters)
    \item Task generalization (MNLI, SQuAD)
    \item Encoder-decoder architectures (T5)
    \item Adequately-powered studies (5+ seeds)
\end{itemize}
